\documentclass{book}
%========================================================
% TABLAS (LaTeX) — ADJETIVOS CLASIFICADOS (MUY AMPLIADAS)
% Incluye: IPA + pronunciación guiada + ejemplo EN + traducción ES
% Recomendación: usar longtable para tablas largas.
%========================================================
% Paquetes sugeridos (agregue en el preámbulo):
\usepackage{longtable}
\usepackage{booktabs}
\usepackage{array}
\usepackage{makecell}
\usepackage[margin=2.5cm]{geometry}
\renewcommand\theadfont{\bfseries}
\setlength\LTpre{6pt}
\setlength\LTpost{6pt}
\newcolumntype{L}[1]{>{\raggedright\arraybackslash}p{#1}}
%========================================================

\begin{document}
	%--------------------------------------------------------
	% 1) Adjetivos descriptivos (Descriptive Adjectives)
	%--------------------------------------------------------
	\begin{longtable}{L{2.8cm} L{3.2cm} L{2.8cm} L{2.6cm} L{6.3cm}}
		\caption{Adjetivos descriptivos (Descriptive Adjectives).}\\
		\toprule
		\thead{Adjetivo} & \thead{Significado} & \thead{IPA} & \thead{Pronunciación} & \thead{Ejemplo de uso (EN / ES)}\\
		\midrule
		\endfirsthead
		\toprule
		\thead{Adjetivo} & \thead{Significado} & \thead{IPA} & \thead{Pronunciación} & \thead{Ejemplo de uso (EN / ES)}\\
		\midrule
		\endhead
		\midrule
		\multicolumn{5}{r}{\small Continúa en la siguiente página.}\\
		\endfoot
		\bottomrule
		\endlastfoot
		
		big & grande & /bɪɡ/ & big &
		\makecell[l]{A \textbf{big} problem needs a solution.\\\emph{Un problema grande necesita una solución.}}\\
		
		small & pequeño & /smɔːl/ & smol &
		\makecell[l]{She lives in a \textbf{small} town.\\\emph{Ella vive en un pueblo pequeño.}}\\
		
		tall & alto & /tɔːl/ & tol &
		\makecell[l]{He is very \textbf{tall} for his age.\\\emph{Él es muy alto para su edad.}}\\
		
		short & bajo / corto & /ʃɔːrt/ & shórt &
		\makecell[l]{The meeting was \textbf{short}.\\\emph{La reunión fue corta.}}\\
		
		long & largo & /lɒŋ/ & long &
		\makecell[l]{It was a \textbf{long} journey.\\\emph{Fue un viaje largo.}}\\
		
		heavy & pesado & /ˈhev.i/ & jé-vi &
		\makecell[l]{This box is \textbf{heavy}.\\\emph{Esta caja es pesada.}}\\
		
		light & ligero & /laɪt/ & lait &
		\makecell[l]{The bag is very \textbf{light}.\\\emph{La bolsa es muy ligera.}}\\
		
		beautiful & hermoso & /ˈbjuː.tɪ.fəl/ & bíu-ti-fol &
		\makecell[l]{It’s a \textbf{beautiful} day.\\\emph{Es un día hermoso.}}\\
		
		ugly & feo & /ˈʌɡ.li/ & ág-li &
		\makecell[l]{That building is \textbf{ugly}.\\\emph{Ese edificio es feo.}}\\
		
		intelligent & inteligente & /ɪnˈtel.ɪ.dʒənt/ & in-té-li-yent &
		\makecell[l]{She is an \textbf{intelligent} student.\\\emph{Ella es una estudiante inteligente.}}\\
		
		clever & listo / astuto & /ˈklev.ər/ & klé-ver &
		\makecell[l]{That was a \textbf{clever} idea.\\\emph{Esa fue una idea astuta.}}\\
		
		strong & fuerte & /strɒŋ/ & strong &
		\makecell[l]{He feels \textbf{strong} today.\\\emph{Él se siente fuerte hoy.}}\\
		
		weak & débil & /wiːk/ & uíik &
		\makecell[l]{The signal is \textbf{weak} here.\\\emph{La señal es débil aquí.}}\\
		
		safe & seguro & /seɪf/ & séif &
		\makecell[l]{This area is \textbf{safe} at night.\\\emph{Esta zona es segura de noche.}}\\
		
		dangerous & peligroso & /ˈdeɪn.dʒər.əs/ & déin-yer-as &
		\makecell[l]{That road is \textbf{dangerous}.\\\emph{Esa carretera es peligrosa.}}\\
		
		noisy & ruidoso & /ˈnɔɪ.zi/ & nói-si &
		\makecell[l]{The classroom is \textbf{noisy}.\\\emph{El aula está ruidosa.}}\\
		
		quiet & silencioso & /ˈkwaɪ.ət/ & kuái-et &
		\makecell[l]{We need a \textbf{quiet} place.\\\emph{Necesitamos un lugar silencioso.}}\\
		
		clean & limpio & /kliːn/ & klíin &
		\makecell[l]{Keep your desk \textbf{clean}.\\\emph{Mantén tu escritorio limpio.}}\\
		
		dirty & sucio & /ˈdɜːr.ti/ & dér-ti &
		\makecell[l]{My shoes are \textbf{dirty}.\\\emph{Mis zapatos están sucios.}}\\
		
		fresh & fresco & /freʃ/ & fresh &
		\makecell[l]{We have \textbf{fresh} information.\\\emph{Tenemos información fresca/reciente.}}\\
		
		stale & pasado (comida) & /steɪl/ & stéil &
		\makecell[l]{The bread is \textbf{stale}.\\\emph{El pan está pasado.}}\\
		
		bright & brillante & /braɪt/ & brait &
		\makecell[l]{It’s a \textbf{bright} room.\\\emph{Es una habitación luminosa.}}\\
		
		dark & oscuro & /dɑːrk/ & dark &
		\makecell[l]{The hallway is \textbf{dark}.\\\emph{El pasillo está oscuro.}}\\
		
		easy & fácil & /ˈiː.zi/ & íi-si &
		\makecell[l]{This exercise is \textbf{easy}.\\\emph{Este ejercicio es fácil.}}\\
		
		hard & difícil / duro & /hɑːrd/ & hard &
		\makecell[l]{The exam was \textbf{hard}.\\\emph{El examen fue difícil.}}\\
		
		simple & simple & /ˈsɪm.pəl/ & sím-pol &
		\makecell[l]{It’s a \textbf{simple} rule.\\\emph{Es una regla simple.}}\\
		
		complex & complejo & /kəmˈpleks/ & kom-pléks &
		\makecell[l]{It’s a \textbf{complex} system.\\\emph{Es un sistema complejo.}}\\
		
	\end{longtable}
	
	%--------------------------------------------------------
	% 2) Adjetivos de opinión (Opinion Adjectives)
	%--------------------------------------------------------
	\begin{longtable}{L{2.8cm} L{3.2cm} L{2.8cm} L{2.6cm} L{6.3cm}}
		\caption{Adjetivos de opinión (Opinion Adjectives).}\\
		\toprule
		\thead{Adjetivo} & \thead{Significado} & \thead{IPA} & \thead{Pronunciación} & \thead{Ejemplo (EN / ES)}\\
		\midrule
		\endfirsthead
		\toprule
		\thead{Adjetivo} & \thead{Significado} & \thead{IPA} & \thead{Pronunciación} & \thead{Ejemplo (EN / ES)}\\
		\midrule
		\endhead
		\midrule
		\multicolumn{5}{r}{\small Continúa en la siguiente página.}\\
		\endfoot
		\bottomrule
		\endlastfoot
		
		good & bueno & /ɡʊd/ & gud &
		\makecell[l]{This is a \textbf{good} example.\\\emph{Este es un buen ejemplo.}}\\
		
		bad & malo & /bæd/ & bad &
		\makecell[l]{It was a \textbf{bad} decision.\\\emph{Fue una mala decisión.}}\\
		
		excellent & excelente & /ˈek.səl.ənt/ & ek-se-lent &
		\makecell[l]{She did an \textbf{excellent} job.\\\emph{Ella hizo un trabajo excelente.}}\\
		
		terrible & terrible & /ˈter.ə.bəl/ & té-re-bol &
		\makecell[l]{The weather is \textbf{terrible}.\\\emph{El clima es terrible.}}\\
		
		amazing & asombroso & /əˈmeɪ.zɪŋ/ & a-méi-zing &
		\makecell[l]{The results are \textbf{amazing}.\\\emph{Los resultados son asombrosos.}}\\
		
		boring & aburrido & /ˈbɔː.rɪŋ/ & bó-ring &
		\makecell[l]{The movie was \textbf{boring}.\\\emph{La película fue aburrida.}}\\
		
		interesting & interesante & /ˈɪn.trə.stɪŋ/ & ín-tres-ting &
		\makecell[l]{It’s an \textbf{interesting} topic.\\\emph{Es un tema interesante.}}\\
		
		useful & útil & /ˈjuːs.fəl/ & yús-fol &
		\makecell[l]{This tool is \textbf{useful}.\\\emph{Esta herramienta es útil.}}\\
		
		useless & inútil & /ˈjuːs.ləs/ & yús-les &
		\makecell[l]{That rule is \textbf{useless}.\\\emph{Esa regla es inútil.}}\\
		
		great & genial / excelente & /ɡreɪt/ & gréit &
		\makecell[l]{That’s a \textbf{great} idea.\\\emph{Esa es una gran idea.}}\\
		
		awful & horrible & /ˈɔː.fəl/ & óo-fol &
		\makecell[l]{The service was \textbf{awful}.\\\emph{El servicio fue horrible.}}\\
		
		fantastic & fantástico & /fænˈtæs.tɪk/ & fan-tás-tik &
		\makecell[l]{We had a \textbf{fantastic} time.\\\emph{La pasamos fantástico.}}\\
		
		wonderful & maravilloso & /ˈwʌn.dər.fəl/ & uán-der-fol &
		\makecell[l]{It’s a \textbf{wonderful} opportunity.\\\emph{Es una oportunidad maravillosa.}}\\
		
		poor & pobre / malo (calidad) & /pʊr/ & pur &
		\makecell[l]{The report is \textbf{poor} in quality.\\\emph{El informe es de mala calidad.}}\\
		
	\end{longtable}
	
	%--------------------------------------------------------
	% 3) Adjetivos de cantidad (Quantitative Adjectives)
	%--------------------------------------------------------
	\begin{longtable}{L{2.8cm} L{3.2cm} L{2.8cm} L{2.6cm} L{6.3cm}}
		\caption{Adjetivos de cantidad (Quantitative Adjectives).}\\
		\toprule
		\thead{Adjetivo} & \thead{Significado} & \thead{IPA} & \thead{Pronunciación} & \thead{Ejemplo (EN / ES)}\\
		\midrule
		\endfirsthead
		\toprule
		\thead{Adjetivo} & \thead{Significado} & \thead{IPA} & \thead{Pronunciación} & \thead{Ejemplo (EN / ES)}\\
		\midrule
		\endhead
		\midrule
		\multicolumn{5}{r}{\small Continúa en la siguiente página.}\\
		\endfoot
		\bottomrule
		\endlastfoot
		
		some & algunos & /sʌm/ & sam &
		\makecell[l]{I need \textbf{some} help.\\\emph{Necesito algo de ayuda.}}\\
		
		any & alguno / ninguno & /ˈen.i/ & é-ni &
		\makecell[l]{Do you have \textbf{any} questions?\\\emph{¿Tienes alguna pregunta?}}\\
		
		many & muchos & /ˈmen.i/ & mé-ni &
		\makecell[l]{There are \textbf{many} students.\\\emph{Hay muchos estudiantes.}}\\
		
		much & mucho & /mʌtʃ/ & mach &
		\makecell[l]{There isn’t \textbf{much} time.\\\emph{No hay mucho tiempo.}}\\
		
		few & pocos & /fjuː/ & fiu &
		\makecell[l]{Only a \textbf{few} people came.\\\emph{Solo vinieron pocas personas.}}\\
		
		little & poco & /ˈlɪt.əl/ & lí-tel &
		\makecell[l]{There is \textbf{little} water left.\\\emph{Queda poca agua.}}\\
		
		several & varios & /ˈsev.ər.əl/ & sé-ve-ral &
		\makecell[l]{She made \textbf{several} attempts.\\\emph{Ella hizo varios intentos.}}\\
		
		enough & suficiente & /ɪˈnʌf/ & i-náf &
		\makecell[l]{We have \textbf{enough} information.\\\emph{Tenemos suficiente información.}}\\
		
		plenty of & abundante & /ˈplen.ti əv/ & plén-ti ov &
		\makecell[l]{There is \textbf{plenty of} food.\\\emph{Hay abundante comida.}}\\
		
		no & ningún & /noʊ/ & nóu &
		\makecell[l]{There is \textbf{no} evidence.\\\emph{No hay evidencia.}}\\
		
		all & todo(s) & /ɔːl/ & ol &
		\makecell[l]{\textbf{All} students must attend.\\\emph{Todos los estudiantes deben asistir.}}\\
		
		both & ambos & /boʊθ/ & bóuz &
		\makecell[l]{\textbf{Both} options are valid.\\\emph{Ambas opciones son válidas.}}\\
		
		each & cada & /iːtʃ/ & íich &
		\makecell[l]{\textbf{Each} student has a code.\\\emph{Cada estudiante tiene un código.}}\\
		
		every & cada (todos) & /ˈev.ri/ & év-ri &
		\makecell[l]{\textbf{Every} day we practice.\\\emph{Cada día practicamos.}}\\
		
	\end{longtable}
	
	%--------------------------------------------------------
	% 4) Adjetivos demostrativos (Demonstrative Adjectives)
	%--------------------------------------------------------
	\begin{longtable}{L{2.8cm} L{3.2cm} L{2.8cm} L{2.6cm} L{6.3cm}}
		\caption{Adjetivos demostrativos (Demonstrative Adjectives).}\\
		\toprule
		\thead{Adjetivo} & \thead{Significado} & \thead{IPA} & \thead{Pronunciación} & \thead{Ejemplo (EN / ES)}\\
		\midrule
		\endfirsthead
		\toprule
		\thead{Adjetivo} & \thead{Significado} & \thead{IPA} & \thead{Pronunciación} & \thead{Ejemplo (EN / ES)}\\
		\midrule
		\endhead
		\bottomrule
		\endfoot
		
		this & este & /ðɪs/ & dhis &
		\makecell[l]{\textbf{This} answer is correct.\\\emph{Esta respuesta es correcta.}}\\
		
		that & ese & /ðæt/ & dhat &
		\makecell[l]{\textbf{That} idea is wrong.\\\emph{Esa idea es incorrecta.}}\\
		
		these & estos & /ðiːz/ & dhiis &
		\makecell[l]{\textbf{These} exercises are easy.\\\emph{Estos ejercicios son fáciles.}}\\
		
		those & esos & /ðoʊz/ & dhóus &
		\makecell[l]{\textbf{Those} books are old.\\\emph{Esos libros son antiguos.}}\\
		
	\end{longtable}
	
	%--------------------------------------------------------
	% 5) Interrogativos (AMPLIADO COMPLETO): todas las question words
	%    en función ADJETIVA (seguida de sustantivo)
	%--------------------------------------------------------
	\begin{longtable}{L{3.3cm} L{3.2cm} L{2.8cm} L{2.6cm} L{5.8cm}}
		\caption{Interrogativos en función adjetiva (Interrogative Adjectives / Question words + noun) — Tabla 5 ampliada.}\\
		\toprule
		\thead{Interrogativo} & \thead{Significado} & \thead{IPA} & \thead{Pronunciación} & \thead{Ejemplo (EN / ES)}\\
		\midrule
		\endfirsthead
		\toprule
		\thead{Interrogativo} & \thead{Significado} & \thead{IPA} & \thead{Pronunciación} & \thead{Ejemplo (EN / ES)}\\
		\midrule
		\endhead
		\midrule
		\multicolumn{5}{r}{\small Continúa en la siguiente página.}\\
		\endfoot
		\bottomrule
		\endlastfoot
		
		which & cuál & /wɪtʃ/ & uích &
		\makecell[l]{\textbf{Which} option do you choose?\\\emph{¿Qué opción eliges?}}\\
		
		what & qué / cuál & /wɒt/ & uót &
		\makecell[l]{\textbf{What} problem are you solving?\\\emph{¿Qué problema estás resolviendo?}}\\
		
		whose & de quién & /huːz/ & júus &
		\makecell[l]{\textbf{Whose} notebook is this?\\\emph{¿De quién es este cuaderno?}}\\
		
		who (as determiner) & cuál persona / qué persona & /huː/ & júu &
		\makecell[l]{\textbf{Who} student is responsible for the report?\\\emph{¿Qué estudiante es responsable del informe?}}\\
		
		where (as determiner) & en qué lugar (de) & /weər/ & uér &
		\makecell[l]{\textbf{Where} campus building is the lab in?\\\emph{¿En qué edificio del campus está el laboratorio?}}\\
		
		when (as determiner) & en qué momento (de) & /wen/ & uén &
		\makecell[l]{\textbf{When} exam date is confirmed?\\\emph{¿Qué fecha de examen está confirmada?}}\\
		
		why (as determiner) & qué razón (de) & /waɪ/ & uái &
		\makecell[l]{\textbf{Why} reason explains the failure?\\\emph{¿Qué razón explica la falla?}}\\
		
		how (as determiner) & de qué manera / qué tan & /haʊ/ & jáu &
		\makecell[l]{\textbf{How} method ensures security?\\\emph{¿Qué método garantiza seguridad?}}\\
		
		what kind of & qué tipo de & /wɒt kaɪnd əv/ & uót kái-nd ov &
		\makecell[l]{\textbf{What kind of} solution do you need?\\\emph{¿Qué tipo de solución necesitas?}}\\
		
		what type of & qué tipo de & /wɒt taɪp əv/ & uót táip ov &
		\makecell[l]{\textbf{What type of} data is required?\\\emph{¿Qué tipo de datos se requieren?}}\\
		
		what sort of & qué clase de & /wɒt sɔːrt əv/ & uót sórt ov &
		\makecell[l]{\textbf{What sort of} evidence supports this?\\\emph{¿Qué clase de evidencia respalda esto?}}\\
		
		how many & cuántos/as & /haʊ ˈmen.i/ & jáu mé-ni &
		\makecell[l]{\textbf{How many} students are present?\\\emph{¿Cuántos estudiantes están presentes?}}\\
		
		how much & cuánto/a (no contable) & /haʊ mʌtʃ/ & jáu mach &
		\makecell[l]{\textbf{How much} time do we have?\\\emph{¿Cuánto tiempo tenemos?}}\\
		
		how long & cuánto tiempo / qué tan largo & /haʊ lɒŋ/ & jáu long &
		\makecell[l]{\textbf{How long} course is this?\\\emph{¿Qué tan largo es este curso?}}\\
		
		how old & qué edad / qué tan antiguo & /haʊ oʊld/ & jáu óuld &
		\makecell[l]{\textbf{How old} system is this?\\\emph{¿Qué tan antiguo es este sistema?}}\\
		
		how far & qué tan lejos & /haʊ fɑːr/ & jáu far &
		\makecell[l]{\textbf{How far} distance is acceptable?\\\emph{¿Qué distancia es aceptable?}}\\
		
		how often & con qué frecuencia & /haʊ ˈɒf.ən/ & jáu ó-fen &
		\makecell[l]{\textbf{How often} meetings are required?\\\emph{¿Con qué frecuencia se requieren reuniones?}}\\
		
		how soon & en cuánto tiempo (pronto) & /haʊ suːn/ & jáu suún &
		\makecell[l]{\textbf{How soon} deadline is realistic?\\\emph{¿Qué plazo es realista?}}\\
		
		how many times & cuántas veces & /haʊ ˈmen.i taɪmz/ & jáu mé-ni táims &
		\makecell[l]{\textbf{How many times} attempts are allowed?\\\emph{¿Cuántos intentos están permitidos?}}\\
		
		which one & cuál (una opción) & /wɪtʃ wʌn/ & uích uán &
		\makecell[l]{\textbf{Which one} answer is correct?\\\emph{¿Cuál respuesta es correcta?}}\\
		
		which type of & qué tipo de & /wɪtʃ taɪp əv/ & uích táip ov &
		\makecell[l]{\textbf{Which type of} algorithm is used?\\\emph{¿Qué tipo de algoritmo se utiliza?}}\\
		
		which level & qué nivel & /wɪtʃ ˈlev.əl/ & uích lé-vel &
		\makecell[l]{\textbf{Which level} course is this?\\\emph{¿De qué nivel es este curso?}}\\
		
	\end{longtable}
	
	%--------------------------------------------------------
	% 6) Adjetivos numerales (Numeral Adjectives) — ampliado
	%--------------------------------------------------------
	\begin{longtable}{L{2.8cm} L{3.2cm} L{2.8cm} L{2.6cm} L{6.3cm}}
		\caption{Adjetivos numerales (Cardinales y ordinales) — ampliado.}\\
		\toprule
		\thead{Forma} & \thead{Significado} & \thead{IPA} & \thead{Pronunciación} & \thead{Ejemplo (EN / ES)}\\
		\midrule
		\endfirsthead
		\toprule
		\thead{Forma} & \thead{Significado} & \thead{IPA} & \thead{Pronunciación} & \thead{Ejemplo (EN / ES)}\\
		\midrule
		\endhead
		\midrule
		\multicolumn{5}{r}{\small Continúa en la siguiente página.}\\
		\endfoot
		\bottomrule
		\endlastfoot
		
		one & uno & /wʌn/ & uán &
		\makecell[l]{I need \textbf{one} answer.\\\emph{Necesito una respuesta.}}\\
		two & dos & /tuː/ & túu &
		\makecell[l]{She has \textbf{two} jobs.\\\emph{Ella tiene dos trabajos.}}\\
		three & tres & /θriː/ & zríi &
		\makecell[l]{They asked \textbf{three} questions.\\\emph{Ellos hicieron tres preguntas.}}\\
		four & cuatro & /fɔːr/ & fór &
		\makecell[l]{The course has \textbf{four} units.\\\emph{El curso tiene cuatro unidades.}}\\
		five & cinco & /faɪv/ & fáiv &
		\makecell[l]{He waited \textbf{five} minutes.\\\emph{Él esperó cinco minutos.}}\\
		six & seis & /sɪks/ & sिक्स (siks) &
		\makecell[l]{We studied \textbf{six} topics.\\\emph{Estudiamos seis temas.}}\\
		seven & siete & /ˈsev.ən/ & sé-ven &
		\makecell[l]{There are \textbf{seven} steps.\\\emph{Hay siete pasos.}}\\
		eight & ocho & /eɪt/ & éit &
		\makecell[l]{She worked \textbf{eight} hours.\\\emph{Ella trabajó ocho horas.}}\\
		nine & nueve & /naɪn/ & náin &
		\makecell[l]{He scored \textbf{nine} points.\\\emph{Él anotó nueve puntos.}}\\
		ten & diez & /ten/ & ten &
		\makecell[l]{There are \textbf{ten} chapters.\\\emph{Hay diez capítulos.}}\\
		eleven & once & /ɪˈlev.ən/ & i-lé-ven &
		\makecell[l]{We have \textbf{eleven} phases.\\\emph{Tenemos once fases.}}\\
		twelve & doce & /twelv/ & tuélv &
		\makecell[l]{\textbf{Twelve} students passed.\\\emph{Doce estudiantes aprobaron.}}\\
		
		first & primero & /fɜːrst/ & férst &
		\makecell[l]{This is the \textbf{first} test.\\\emph{Este es el primer examen.}}\\
		second & segundo & /ˈsek.ənd/ & sé-cond &
		\makecell[l]{She came \textbf{second}.\\\emph{Ella llegó segunda.}}\\
		third & tercero & /θɜːrd/ & zérd &
		\makecell[l]{It’s the \textbf{third} attempt.\\\emph{Es el tercer intento.}}\\
		fourth & cuarto & /fɔːrθ/ & fórz &
		\makecell[l]{The \textbf{fourth} chapter is long.\\\emph{El cuarto capítulo es largo.}}\\
		fifth & quinto & /fɪfθ/ & fifz &
		\makecell[l]{The \textbf{fifth} task is optional.\\\emph{La quinta tarea es opcional.}}\\
		last & último & /læst/ & last &
		\makecell[l]{This is the \textbf{last} example.\\\emph{Este es el último ejemplo.}}\\
		
	\end{longtable}
	
	%--------------------------------------------------------
	% 7) Adjetivos de estado/condición (State / Condition)
	%--------------------------------------------------------
	\begin{longtable}{L{2.8cm} L{3.2cm} L{2.8cm} L{2.6cm} L{6.3cm}}
		\caption{Adjetivos de estado o condición (State / Condition Adjectives).}\\
		\toprule
		\thead{Adjetivo} & \thead{Significado} & \thead{IPA} & \thead{Pronunciación} & \thead{Ejemplo (EN / ES)}\\
		\midrule
		\endfirsthead
		\toprule
		\thead{Adjetivo} & \thead{Significado} & \thead{IPA} & \thead{Pronunciación} & \thead{Ejemplo (EN / ES)}\\
		\midrule
		\endhead
		\midrule
		\multicolumn{5}{r}{\small Continúa en la siguiente página.}\\
		\endfoot
		\bottomrule
		\endlastfoot
		
		happy & feliz & /ˈhæp.i/ & já-pi &
		\makecell[l]{I feel \textbf{happy} today.\\\emph{Me siento feliz hoy.}}\\
		sad & triste & /sæd/ & sad &
		\makecell[l]{He looks \textbf{sad}.\\\emph{Él se ve triste.}}\\
		tired & cansado & /ˈtaɪərd/ & tái-erd &
		\makecell[l]{She is \textbf{tired} after work.\\\emph{Ella está cansada después del trabajo.}}\\
		angry & enojado & /ˈæŋ.ɡri/ & án-gri &
		\makecell[l]{The client is \textbf{angry}.\\\emph{El cliente está enojado.}}\\
		sick & enfermo & /sɪk/ & sik &
		\makecell[l]{He is \textbf{sick} today.\\\emph{Él está enfermo hoy.}}\\
		healthy & saludable & /ˈhel.θi/ & jél-thi &
		\makecell[l]{She eats \textbf{healthy} food.\\\emph{Ella come comida saludable.}}\\
		ready & listo & /ˈred.i/ & ré-di &
		\makecell[l]{We are \textbf{ready} now.\\\emph{Estamos listos ahora.}}\\
		afraid & asustado & /əˈfreɪd/ & a-fréid &
		\makecell[l]{She is \textbf{afraid} of dogs.\\\emph{Ella tiene miedo a los perros.}}\\
		calm & tranquilo & /kɑːm/ & kam &
		\makecell[l]{He stayed \textbf{calm}.\\\emph{Él se mantuvo tranquilo.}}\\
		busy & ocupado & /ˈbɪz.i/ & bí-si &
		\makecell[l]{I’m \textbf{busy} right now.\\\emph{Estoy ocupado ahora mismo.}}\\
		free & libre & /friː/ & fríi &
		\makecell[l]{Are you \textbf{free} tomorrow?\\\emph{¿Estás libre mañana?}}\\
		
	\end{longtable}
	
	%--------------------------------------------------------
	% 8) Adjetivos de edad (Age Adjectives)
	%--------------------------------------------------------
	\begin{longtable}{L{2.8cm} L{3.2cm} L{2.8cm} L{2.6cm} L{6.3cm}}
		\caption{Adjetivos de edad (Age Adjectives).}\\
		\toprule
		\thead{Adjetivo} & \thead{Significado} & \thead{IPA} & \thead{Pronunciación} & \thead{Ejemplo (EN / ES)}\\
		\midrule
		\endfirsthead
		\toprule
		\thead{Adjetivo} & \thead{Significado} & \thead{IPA} & \thead{Pronunciación} & \thead{Ejemplo (EN / ES)}\\
		\midrule
		\endhead
		\bottomrule
		\endfoot
		
		young & joven & /jʌŋ/ & yáng &
		\makecell[l]{A \textbf{young} teacher arrived.\\\emph{Llegó un profesor joven.}}\\
		old & viejo / antiguo & /oʊld/ & óuld &
		\makecell[l]{The system is \textbf{old}.\\\emph{El sistema es antiguo.}}\\
		new & nuevo & /njuː/ & niú &
		\makecell[l]{It’s a \textbf{new} version.\\\emph{Es una versión nueva.}}\\
		ancient & antiguo (muy viejo) & /ˈeɪn.ʃənt/ & éin-shent &
		\makecell[l]{An \textbf{ancient} culture existed here.\\\emph{Una cultura antigua existió aquí.}}\\
		modern & moderno & /ˈmɒd.ərn/ & mó-dern &
		\makecell[l]{A \textbf{modern} solution is needed.\\\emph{Se necesita una solución moderna.}}\\
		
	\end{longtable}
	
	%--------------------------------------------------------
	% 9) Adjetivos de forma (Shape Adjectives)
	%--------------------------------------------------------
	\begin{longtable}{L{2.8cm} L{3.2cm} L{2.8cm} L{2.6cm} L{6.3cm}}
		\caption{Adjetivos de forma (Shape Adjectives).}\\
		\toprule
		\thead{Adjetivo} & \thead{Significado} & \thead{IPA} & \thead{Pronunciación} & \thead{Ejemplo (EN / ES)}\\
		\midrule
		\endfirsthead
		\toprule
		\thead{Adjetivo} & \thead{Significado} & \thead{IPA} & \thead{Pronunciación} & \thead{Ejemplo (EN / ES)}\\
		\midrule
		\endhead
		\bottomrule
		\endfoot
		
		round & redondo & /raʊnd/ & ráund &
		\makecell[l]{A \textbf{round} table fits here.\\\emph{Una mesa redonda encaja aquí.}}\\
		square & cuadrado & /skweər/ & skuér &
		\makecell[l]{The room is \textbf{square}.\\\emph{La habitación es cuadrada.}}\\
		flat & plano & /flæt/ & flat &
		\makecell[l]{The surface is \textbf{flat}.\\\emph{La superficie es plana.}}\\
		straight & recto & /streɪt/ & stréit &
		\makecell[l]{Draw a \textbf{straight} line.\\\emph{Dibuja una línea recta.}}\\
		curved & curvo & /kɜːrvd/ & kérvd &
		\makecell[l]{The road is \textbf{curved}.\\\emph{La carretera es curva.}}\\
		
	\end{longtable}
	
	%--------------------------------------------------------
	% 10) Adjetivos de color (Color Adjectives)
	%--------------------------------------------------------
	\begin{longtable}{L{2.8cm} L{3.2cm} L{2.8cm} L{2.6cm} L{6.3cm}}
		\caption{Adjetivos de color (Color Adjectives).}\\
		\toprule
		\thead{Adjetivo} & \thead{Significado} & \thead{IPA} & \thead{Pronunciación} & \thead{Ejemplo (EN / ES)}\\
		\midrule
		\endfirsthead
		\toprule
		\thead{Adjetivo} & \thead{Significado} & \thead{IPA} & \thead{Pronunciación} & \thead{Ejemplo (EN / ES)}\\
		\midrule
		\endhead
		\bottomrule
		\endfoot
		
		red & rojo & /red/ & red &
		\makecell[l]{A \textbf{red} light means stop.\\\emph{Una luz roja significa detenerse.}}\\
		blue & azul & /bluː/ & blúu &
		\makecell[l]{She bought a \textbf{blue} dress.\\\emph{Ella compró un vestido azul.}}\\
		green & verde & /ɡriːn/ & gríin &
		\makecell[l]{A \textbf{green} area is nearby.\\\emph{Hay una zona verde cerca.}}\\
		black & negro & /blæk/ & blak &
		\makecell[l]{He wears \textbf{black} shoes.\\\emph{Él usa zapatos negros.}}\\
		white & blanco & /waɪt/ & uait &
		\makecell[l]{The walls are \textbf{white}.\\\emph{Las paredes son blancas.}}\\
		yellow & amarillo & /ˈjel.oʊ/ & yé-lo &
		\makecell[l]{A \textbf{yellow} sign appeared.\\\emph{Apareció una señal amarilla.}}\\
		brown & marrón & /braʊn/ & ब्राउन (braun) &
		\makecell[l]{A \textbf{brown} bag is missing.\\\emph{Falta una bolsa marrón.}}\\
		
	\end{longtable}
	
	%--------------------------------------------------------
	% 11) Adjetivos de origen (Origin Adjectives)
	%--------------------------------------------------------
	\begin{longtable}{L{2.8cm} L{3.2cm} L{2.8cm} L{2.6cm} L{6.3cm}}
		\caption{Adjetivos de origen (Origin Adjectives).}\\
		\toprule
		\thead{Adjetivo} & \thead{Significado} & \thead{IPA} & \thead{Pronunciación} & \thead{Ejemplo (EN / ES)}\\
		\midrule
		\endfirsthead
		\toprule
		\thead{Adjetivo} & \thead{Significado} & \thead{IPA} & \thead{Pronunciación} & \thead{Ejemplo (EN / ES)}\\
		\midrule
		\endhead
		\bottomrule
		\endfoot
		
		Ecuadorian & ecuatoriano & /ˌek.wəˈdɔːr.i.ən/ & ek-wa-dór-ian &
		\makecell[l]{\textbf{Ecuadorian} food is popular.\\\emph{La comida ecuatoriana es popular.}}\\
		American & estadounidense & /əˈmer.ɪ.kən/ & a-mé-ri-kan &
		\makecell[l]{An \textbf{American} company invested.\\\emph{Una empresa estadounidense invirtió.}}\\
		Spanish & español & /ˈspæn.ɪʃ/ & स्पै-nish (spá-nish) &
		\makecell[l]{She studies \textbf{Spanish} literature.\\\emph{Ella estudia literatura española.}}\\
		Chinese & chino & /ˌtʃaɪˈniːz/ & chai-níiz &
		\makecell[l]{\textbf{Chinese} technology is growing.\\\emph{La tecnología china está creciendo.}}\\
		European & europeo & /ˌjʊə.rəˈpiː.ən/ & yú-ra-pí-an &
		\makecell[l]{\textbf{European} standards apply here.\\\emph{Las normas europeas se aplican aquí.}}\\
		
	\end{longtable}
	
	%--------------------------------------------------------
	% 12) Adjetivos de material (Material Adjectives)
	%--------------------------------------------------------
	\begin{longtable}{L{2.8cm} L{3.2cm} L{2.8cm} L{2.6cm} L{6.3cm}}
		\caption{Adjetivos de material (Material Adjectives).}\\
		\toprule
		\thead{Adjetivo} & \thead{Significado} & \thead{IPA} & \thead{Pronunciación} & \thead{Ejemplo (EN / ES)}\\
		\midrule
		\endfirsthead
		\toprule
		\thead{Adjetivo} & \thead{Significado} & \thead{IPA} & \thead{Pronunciación} & \thead{Ejemplo (EN / ES)}\\
		\midrule
		\endhead
		\bottomrule
		\endfoot
		
		wooden & de madera & /ˈwʊd.ən/ & वुड-न (wúd-en) &
		\makecell[l]{A \textbf{wooden} chair broke.\\\emph{Se rompió una silla de madera.}}\\
		metal & metálico & /ˈmet.əl/ & mé-tel &
		\makecell[l]{A \textbf{metal} structure stands.\\\emph{Se mantiene una estructura metálica.}}\\
		plastic & de plástico & /ˈplæs.tɪk/ & प्लás-tik (plás-tik) &
		\makecell[l]{A \textbf{plastic} bottle fell.\\\emph{Se cayó una botella de plástico.}}\\
		glass & de vidrio & /ɡlæs/ & glas &
		\makecell[l]{A \textbf{glass} door opened.\\\emph{Se abrió una puerta de vidrio.}}\\
		cotton & de algodón & /ˈkɒt.ən/ & kó-ton &
		\makecell[l]{A \textbf{cotton} shirt is soft.\\\emph{Una camisa de algodón es suave.}}\\
		leather & de cuero & /ˈleð.ər/ & lé-ther &
		\makecell[l]{A \textbf{leather} jacket is expensive.\\\emph{Una chaqueta de cuero es cara.}}\\
		
	\end{longtable}
	
	%--------------------------------------------------------
	% 13) Adjetivos de propósito (Purpose Adjectives)
	%--------------------------------------------------------
	\begin{longtable}{L{2.8cm} L{3.2cm} L{2.8cm} L{2.6cm} L{6.3cm}}
		\caption{Adjetivos de propósito (Purpose Adjectives).}\\
		\toprule
		\thead{Adjetivo} & \thead{Significado} & \thead{IPA} & \thead{Pronunciación} & \thead{Ejemplo (EN / ES)}\\
		\midrule
		\endfirsthead
		\toprule
		\thead{Adjetivo} & \thead{Significado} & \thead{IPA} & \thead{Pronunciación} & \thead{Ejemplo (EN / ES)}\\
		\midrule
		\endhead
		\bottomrule
		\endfoot
		
		sleeping & para dormir & /ˈsliː.pɪŋ/ & slíi-ping &
		\makecell[l]{A \textbf{sleeping} bag is useful.\\\emph{Una bolsa para dormir es útil.}}\\
		working & de trabajo & /ˈwɜːr.kɪŋ/ & uér-king &
		\makecell[l]{\textbf{Working} hours changed.\\\emph{Las horas de trabajo cambiaron.}}\\
		cooking & para cocinar & /ˈkʊk.ɪŋ/ & kú-king &
		\makecell[l]{\textbf{Cooking} tools are ready.\\\emph{Las herramientas de cocina están listas.}}\\
		writing & para escribir & /ˈraɪ.tɪŋ/ & ráy-ting &
		\makecell[l]{A \textbf{writing} desk helps.\\\emph{Un escritorio para escribir ayuda.}}\\
		reading & para leer & /ˈriː.dɪŋ/ & ríi-ding &
		\makecell[l]{A \textbf{reading} lamp is on.\\\emph{Una lámpara de lectura está encendida.}}\\
		
	\end{longtable}
	
	%--------------------------------------------------------
	% 14) Comparativos y superlativos — AÚN MÁS AMPLIADO
	% Incluye: más filas + IPA base/comp/sup + 3 oraciones (EN/ES) por fila
	%--------------------------------------------------------
	\begin{longtable}{L{2.2cm} L{2.6cm} L{2.8cm} L{2.6cm} L{2.8cm} L{3.0cm} L{7.0cm}}
		\caption{Comparativos y superlativos (Tabla 14 ampliada: IPA para base/comparativo/superlativo y 3 oraciones por fila).}\\
		\toprule
		\thead{Base} &
		\thead{Comparativo} &
		\thead{Superlativo} &
		\thead{IPA (Base)} &
		\thead{IPA (Comp.)} &
		\thead{IPA (Sup.)} &
		\thead{Oraciones (Base / Comp. / Sup.) (EN / ES)}\\
		\midrule
		\endfirsthead
		\toprule
		\thead{Base} &
		\thead{Comparativo} &
		\thead{Superlativo} &
		\thead{IPA (Base)} &
		\thead{IPA (Comp.)} &
		\thead{IPA (Sup.)} &
		\thead{Oraciones (Base / Comp. / Sup.) (EN / ES)}\\
		\midrule
		\endhead
		\midrule
		\multicolumn{7}{r}{\small Continúa en la siguiente página.}\\
		\endfoot
		\bottomrule
		\endlastfoot
		
		%------------------- (ya existentes) --------------------
		big & bigger & biggest & /bɪɡ/ & /ˈbɪɡ.ər/ & /ˈbɪɡ.ɪst/ &
		\makecell[l]{\textbf{Base:} This is a \textbf{big} project.\\\emph{Este es un proyecto grande.}\\
			\textbf{Comp.:} This project is \textbf{bigger} than the last one.\\\emph{Este proyecto es más grande que el anterior.}\\
			\textbf{Sup.:} This is the \textbf{biggest} project this year.\\\emph{Este es el proyecto más grande de este año.}}\\
		
		fast & faster & fastest & /fæst/ & /ˈfæs.tər/ & /ˈfæs.tɪst/ &
		\makecell[l]{\textbf{Base:} He is \textbf{fast}.\\\emph{Él es rápido.}\\
			\textbf{Comp.:} He is \textbf{faster} now.\\\emph{Él es más rápido ahora.}\\
			\textbf{Sup.:} He is the \textbf{fastest} runner in the class.\\\emph{Él es el corredor más rápido de la clase.}}\\
		
		good & better & best & /ɡʊd/ & /ˈbet.ər/ & /best/ &
		\makecell[l]{\textbf{Base:} This is a \textbf{good} idea.\\\emph{Esta es una buena idea.}\\
			\textbf{Comp.:} This idea is \textbf{better} than yours.\\\emph{Esta idea es mejor que la tuya.}\\
			\textbf{Sup.:} This is the \textbf{best} option.\\\emph{Esta es la mejor opción.}}\\
		
		bad & worse & worst & /bæd/ & /wɜːrs/ & /wɜːrst/ &
		\makecell[l]{\textbf{Base:} That is a \textbf{bad} plan.\\\emph{Ese es un mal plan.}\\
			\textbf{Comp.:} This plan is \textbf{worse} than the first one.\\\emph{Este plan es peor que el primero.}\\
			\textbf{Sup.:} That was the \textbf{worst} mistake.\\\emph{Ese fue el peor error.}}\\
		
		easy & easier & easiest & /ˈiː.zi/ & /ˈiː.zi.ər/ & /ˈiː.zi.ɪst/ &
		\makecell[l]{\textbf{Base:} This exercise is \textbf{easy}.\\\emph{Este ejercicio es fácil.}\\
			\textbf{Comp.:} This exercise is \textbf{easier} than the previous one.\\\emph{Este ejercicio es más fácil que el anterior.}\\
			\textbf{Sup.:} This is the \textbf{easiest} question.\\\emph{Esta es la pregunta más fácil.}}\\
		
		hard & harder & hardest & /hɑːrd/ & /ˈhɑːr.dər/ & /ˈhɑːr.dɪst/ &
		\makecell[l]{\textbf{Base:} The test is \textbf{hard}.\\\emph{El examen es difícil.}\\
			\textbf{Comp.:} Today’s test is \textbf{harder} than yesterday’s.\\\emph{El examen de hoy es más difícil que el de ayer.}\\
			\textbf{Sup.:} This is the \textbf{hardest} topic.\\\emph{Este es el tema más difícil.}}\\
		
		high & higher & highest & /haɪ/ & /ˈhaɪ.ər/ & /ˈhaɪ.ɪst/ &
		\makecell[l]{\textbf{Base:} The price is \textbf{high}.\\\emph{El precio es alto.}\\
			\textbf{Comp.:} The price is \textbf{higher} this month.\\\emph{El precio es más alto este mes.}\\
			\textbf{Sup.:} This is the \textbf{highest} price of the year.\\\emph{Este es el precio más alto del año.}}\\
		
		low & lower & lowest & /loʊ/ & /ˈloʊ.ər/ & /ˈloʊ.ɪst/ &
		\makecell[l]{\textbf{Base:} The risk is \textbf{low}.\\\emph{El riesgo es bajo.}\\
			\textbf{Comp.:} The risk is \textbf{lower} with this method.\\\emph{El riesgo es más bajo con este método.}\\
			\textbf{Sup.:} This is the \textbf{lowest} risk option.\\\emph{Esta es la opción de menor riesgo.}}\\
		
		near & nearer & nearest & /nɪər/ & /ˈnɪə.rər/ & /ˈnɪə.rɪst/ &
		\makecell[l]{\textbf{Base:} The office is \textbf{near}.\\\emph{La oficina está cerca.}\\
			\textbf{Comp.:} This office is \textbf{nearer} than the other one.\\\emph{Esta oficina está más cerca que la otra.}\\
			\textbf{Sup.:} This is the \textbf{nearest} office.\\\emph{Esta es la oficina más cercana.}}\\
		
		far & farther / further & farthest / furthest & /fɑːr/ & /ˈfɑːr.ðər/ /ˈfɜːr.ðər/ & /ˈfɑːr.ðɪst/ /ˈfɜːr.ðɪst/ &
		\makecell[l]{\textbf{Base:} The campus is \textbf{far}.\\\emph{El campus está lejos.}\\
			\textbf{Comp.:} The library is \textbf{farther} than the lab.\\\emph{La biblioteca está más lejos que el laboratorio.}\\
			\textbf{Sup.:} That is the \textbf{farthest} building.\\\emph{Ese es el edificio más lejano.}}\\
		
		%------------------- (NUEVAS FILAS) --------------------
		small & smaller & smallest & /smɔːl/ & /ˈsmɔːl.ər/ & /ˈsmɔːl.ɪst/ &
		\makecell[l]{\textbf{Base:} This is a \textbf{small} room.\\\emph{Esta es una habitación pequeña.}\\
			\textbf{Comp.:} This room is \textbf{smaller} than mine.\\\emph{Esta habitación es más pequeña que la mía.}\\
			\textbf{Sup.:} This is the \textbf{smallest} room in the house.\\\emph{Esta es la habitación más pequeña de la casa.}}\\
		
		tall & taller & tallest & /tɔːl/ & /ˈtɔːl.ər/ & /ˈtɔːl.ɪst/ &
		\makecell[l]{\textbf{Base:} He is \textbf{tall}.\\\emph{Él es alto.}\\
			\textbf{Comp.:} He is \textbf{taller} than his brother.\\\emph{Él es más alto que su hermano.}\\
			\textbf{Sup.:} He is the \textbf{tallest} student here.\\\emph{Él es el estudiante más alto aquí.}}\\
		
		short & shorter & shortest & /ʃɔːrt/ & /ˈʃɔːr.tər/ & /ˈʃɔːr.tɪst/ &
		\makecell[l]{\textbf{Base:} The lesson is \textbf{short}.\\\emph{La lección es corta.}\\
			\textbf{Comp.:} Today’s lesson is \textbf{shorter}.\\\emph{La lección de hoy es más corta.}\\
			\textbf{Sup.:} This is the \textbf{shortest} video.\\\emph{Este es el video más corto.}}\\
		
		long & longer & longest & /lɒŋ/ & /ˈlɒŋ.ɡər/ & /ˈlɒŋ.ɡɪst/ &
		\makecell[l]{\textbf{Base:} It’s a \textbf{long} story.\\\emph{Es una historia larga.}\\
			\textbf{Comp.:} This story is \textbf{longer} than that one.\\\emph{Esta historia es más larga que esa.}\\
			\textbf{Sup.:} This is the \textbf{longest} chapter.\\\emph{Este es el capítulo más largo.}}\\
		
		new & newer & newest & /njuː/ & /ˈnjuː.ər/ & /ˈnjuː.ɪst/ &
		\makecell[l]{\textbf{Base:} I bought a \textbf{new} laptop.\\\emph{Compré una laptop nueva.}\\
			\textbf{Comp.:} This laptop is \textbf{newer} than mine.\\\emph{Esta laptop es más nueva que la mía.}\\
			\textbf{Sup.:} This is the \textbf{newest} model.\\\emph{Este es el modelo más nuevo.}}\\
		
		old & older & oldest & /oʊld/ & /ˈoʊl.dər/ & /ˈoʊl.dɪst/ &
		\makecell[l]{\textbf{Base:} That building is \textbf{old}.\\\emph{Ese edificio es antiguo.}\\
			\textbf{Comp.:} This building is \textbf{older} than the library.\\\emph{Este edificio es más antiguo que la biblioteca.}\\
			\textbf{Sup.:} It is the \textbf{oldest} school in town.\\\emph{Es la escuela más antigua del pueblo.}}\\
		
		young & younger & youngest & /jʌŋ/ & /ˈjʌŋ.ɡər/ & /ˈjʌŋ.ɡɪst/ &
		\makecell[l]{\textbf{Base:} She is \textbf{young}.\\\emph{Ella es joven.}\\
			\textbf{Comp.:} She is \textbf{younger} than me.\\\emph{Ella es menor que yo.}\\
			\textbf{Sup.:} He is the \textbf{youngest} member.\\\emph{Él es el miembro más joven.}}\\
		
		rich & richer & richest & /rɪtʃ/ & /ˈrɪtʃ.ər/ & /ˈrɪtʃ.ɪst/ &
		\makecell[l]{\textbf{Base:} That city is \textbf{rich}.\\\emph{Esa ciudad es rica.}\\
			\textbf{Comp.:} This region is \textbf{richer} in resources.\\\emph{Esta región es más rica en recursos.}\\
			\textbf{Sup.:} It is the \textbf{richest} neighborhood.\\\emph{Es el barrio más rico.}}\\
		
		poor & poorer & poorest & /pʊr/ & /ˈpʊr.ər/ & /ˈpʊr.ɪst/ &
		\makecell[l]{\textbf{Base:} The country is \textbf{poor}.\\\emph{El país es pobre.}\\
			\textbf{Comp.:} This area is \textbf{poorer} than the center.\\\emph{Esta zona es más pobre que el centro.}\\
			\textbf{Sup.:} It is the \textbf{poorest} district.\\\emph{Es el distrito más pobre.}}\\
		
		safe & safer & safest & /seɪf/ & /ˈseɪ.fər/ & /ˈseɪ.fɪst/ &
		\makecell[l]{\textbf{Base:} This road is \textbf{safe}.\\\emph{Esta carretera es segura.}\\
			\textbf{Comp.:} This route is \textbf{safer} at night.\\\emph{Esta ruta es más segura de noche.}\\
			\textbf{Sup.:} This is the \textbf{safest} option.\\\emph{Esta es la opción más segura.}}\\
		
		cheap & cheaper & cheapest & /tʃiːp/ & /ˈtʃiː.pər/ & /ˈtʃiː.pɪst/ &
		\makecell[l]{\textbf{Base:} This phone is \textbf{cheap}.\\\emph{Este teléfono es barato.}\\
			\textbf{Comp.:} This phone is \textbf{cheaper} online.\\\emph{Este teléfono es más barato en línea.}\\
			\textbf{Sup.:} This is the \textbf{cheapest} store.\\\emph{Esta es la tienda más barata.}}\\
		
		expensive & more expensive & most expensive & /ɪkˈspen.sɪv/ & /mɔːr ɪkˈspen.sɪv/ & /moʊst ɪkˈspen.sɪv/ &
		\makecell[l]{\textbf{Base:} This laptop is \textbf{expensive}.\\\emph{Esta laptop es cara.}\\
			\textbf{Comp.:} This laptop is \textbf{more expensive} than that one.\\\emph{Esta laptop es más cara que esa.}\\
			\textbf{Sup.:} It is the \textbf{most expensive} model.\\\emph{Es el modelo más caro.}}\\
		
		important & more important & most important & /ɪmˈpɔːr.tənt/ & /mɔːr ɪmˈpɔːr.tənt/ & /moʊst ɪmˈpɔːr.tənt/ &
		\makecell[l]{\textbf{Base:} This rule is \textbf{important}.\\\emph{Esta regla es importante.}\\
			\textbf{Comp.:} Security is \textbf{more important} than speed.\\\emph{La seguridad es más importante que la velocidad.}\\
			\textbf{Sup.:} This is the \textbf{most important} part.\\\emph{Esta es la parte más importante.}}\\
		
		interesting & more interesting & most interesting & /ˈɪn.trə.stɪŋ/ & /mɔːr ˈɪn.trə.stɪŋ/ & /moʊst ˈɪn.trə.stɪŋ/ &
		\makecell[l]{\textbf{Base:} The topic is \textbf{interesting}.\\\emph{El tema es interesante.}\\
			\textbf{Comp.:} This topic is \textbf{more interesting} than math.\\\emph{Este tema es más interesante que matemáticas.}\\
			\textbf{Sup.:} This is the \textbf{most interesting} chapter.\\\emph{Este es el capítulo más interesante.}}\\
		
		useful & more useful & most useful & /ˈjuːs.fəl/ & /mɔːr ˈjuːs.fəl/ & /moʊst ˈjuːs.fəl/ &
		\makecell[l]{\textbf{Base:} This tool is \textbf{useful}.\\\emph{Esta herramienta es útil.}\\
			\textbf{Comp.:} This tool is \textbf{more useful} for beginners.\\\emph{Esta herramienta es más útil para principiantes.}\\
			\textbf{Sup.:} This is the \textbf{most useful} feature.\\\emph{Esta es la función más útil.}}\\
		
		simple & simpler & simplest & /ˈsɪm.pəl/ & /ˈsɪm.plər/ & /ˈsɪm.plɪst/ &
		\makecell[l]{\textbf{Base:} The explanation is \textbf{simple}.\\\emph{La explicación es simple.}\\
			\textbf{Comp.:} This example is \textbf{simpler}.\\\emph{Este ejemplo es más simple.}\\
			\textbf{Sup.:} This is the \textbf{simplest} method.\\\emph{Este es el método más simple.}}\\
		
		complex & more complex & most complex & /kəmˈpleks/ & /mɔːr kəmˈpleks/ & /moʊst kəmˈpleks/ &
		\makecell[l]{\textbf{Base:} The system is \textbf{complex}.\\\emph{El sistema es complejo.}\\
			\textbf{Comp.:} This system is \textbf{more complex} than the prototype.\\\emph{Este sistema es más complejo que el prototipo.}\\
			\textbf{Sup.:} This is the \textbf{most complex} module.\\\emph{Este es el módulo más complejo.}}\\
		
		hot & hotter & hottest & /hɒt/ & /ˈhɒt.ər/ & /ˈhɒt.ɪst/ &
		\makecell[l]{\textbf{Base:} The coffee is \textbf{hot}.\\\emph{El café está caliente.}\\
			\textbf{Comp.:} Today is \textbf{hotter} than yesterday.\\\emph{Hoy hace más calor que ayer.}\\
			\textbf{Sup.:} August is the \textbf{hottest} month.\\\emph{Agosto es el mes más caluroso.}}\\
		
		cold & colder & coldest & /koʊld/ & /ˈkoʊl.dər/ & /ˈkoʊl.dɪst/ &
		\makecell[l]{\textbf{Base:} The water is \textbf{cold}.\\\emph{El agua está fría.}\\
			\textbf{Comp.:} This room is \textbf{colder}.\\\emph{Esta habitación está más fría.}\\
			\textbf{Sup.:} January is the \textbf{coldest} month.\\\emph{Enero es el mes más frío.}}\\
		
	\end{longtable}
	
\end{document}